% source: RKF/theorum/04_zero_cut_t03_endpoint_completion.md
\hypertarget{theorum-04-zero-cut-t03-endpoint-completion}{%
\chapter{Theorum 04: Zero-Cut T03 Endpoint Completion}\label{theorum-04-zero-cut-t03-endpoint-completion}}

\hypertarget{source-provenance}{%
\section{Source provenance}\label{source-provenance}}

\begin{verbatim}
source repository: Parveen117/MP
primary PR:        #254 Close the five-dimensional completed-Weil endpoint by explicit zero-cut T03 completion
support PRs:       #242, #247, #255, #256
main objects:      R_Sigma, P_src, S_(0,-), W_Sigma, A3, K_0
\end{verbatim}

\hypertarget{purpose}{%
\section{Purpose}\label{purpose}}

This theorem package records the endpoint completion of the five-dimensional completed-Weil packet after the seam-integer and single-sector reductions have fixed the terminal object.

It should not be confused with the earlier exploratory finite eta sequence. The T03 theorem is the endpoint source-readout theorem.

\hypertarget{native-t03-specialization}{%
\section{Native T03 specialization}\label{native-t03-specialization}}

The zero-cut theorem fixes, before any lift is chosen:

\begin{verbatim}
projected native observer;
completed-Weil target;
actual target-relevant blind quotient.
\end{verbatim}

The two spectral inputs are:

\begin{verbatim}
Spectral I:
  exact five-dimensional defect;
  minimal five-observer repair.

Spectral II:
  target-relative lift count;
  positive uniform repair margin;
  a posteriori response-defect control;
  matched visible/blind/cross certification;
  strong completion.
\end{verbatim}

\hypertarget{six-block-gram-decomposition}{%
\section{Six-block Gram decomposition}\label{six-block-gram-decomposition}}

For transported finite source maps \texttt{Y\_n} and \texttt{T\_n}, the five-label Gram defect is decomposed into

\begin{verbatim}
visible compression
+ five-dimensional blind response
+ both cross blocks
+ omitted prime/Gamma tails
+ quadrature/window/solve/rounding errors
+ native/Feshbach/prime-torus chart transport.
\end{verbatim}

The theorem records that the six directed bounds have total error tending to zero.

\hypertarget{invariant-source-gram-identity}{%
\section{Invariant source-Gram identity}\label{invariant-source-gram-identity}}

Strong completion proves the invariant source-Gram identity

\begin{verbatim}
R_Sigma^* P_src R_Sigma
  = (S_(0,-)^(dagger/2) W_Sigma)^*
    (S_(0,-)^(dagger/2) W_Sigma).
\end{verbatim}

Literal equality of ambient source maps is not assumed. Equal Grams give a canonical partial-isometry identification of the final ranges. The endpoint sign needs only the invariant Gram identity.

\hypertarget{projection-defect-positivity}{%
\section{Projection-defect positivity}\label{projection-defect-positivity}}

The independently constructed total-response Gram is

\begin{verbatim}
R_Sigma^* R_Sigma = diag(1,A3).
\end{verbatim}

Therefore the endpoint projection defect is

\begin{verbatim}
K_0 = R_Sigma^*(I-P_src)R_Sigma >= 0.
\end{verbatim}

The endpoint Schur equivalence then transfers this projection-defect positivity to the completed-Weil endpoint sign in the declared normalization.

\hypertarget{terminal-membrane}{%
\section{Terminal membrane}\label{terminal-membrane}}

The theorem uses the classical completed explicit-formula identification and Weil criterion as the final classical membrane.

\begin{verbatim}
native endpoint Schur sign
+ classical completed explicit formula / Weil criterion
-> RH in the declared normalization.
\end{verbatim}

This folder records the native theorem and its stated terminal interface. It does not pretend that the classical membrane was independently rederived inside this capsule.

\hypertarget{claim-boundary}{%
\section{Claim boundary}\label{claim-boundary}}

\begin{verbatim}
zero-cut quotient                              FIXED
five-label lift                                SUPPLIED BY SPECTRAL I-II
six-block Gram decomposition                   SUPPLIED
vanishing total error                           SUPPLIED
orientation-free source readout                 SUPPLIED
K_0 = R_Sigma^*(I-P_src)R_Sigma >= 0            PROVED IN SOURCE THEOREM
endpoint Schur closure                          PROVED IN SOURCE THEOREM
classical Weil terminal membrane                CITED/CONSUMED WHERE STATED
\end{verbatim}

\hypertarget{use-in-the-framework}{%
\section{Use in the framework}\label{use-in-the-framework}}

Consume this theorem as the endpoint closure adapter:

\begin{verbatim}
single-sector / five-label source packet
-> invariant source-Gram identity
-> K_0 >= 0
-> endpoint completed-Weil sign
-> classical membrane if invoked
\end{verbatim}

Do not reopen M3, path equality, or seam integer while using T03. They are upstream adapters already transferred into this folder.

\begin{flushright}\small\texttt{RKF/theorum/04\_zero\_cut\_t03\_endpoint\_completion.md}\end{flushright}
